\section{Template Features}

This document serves as a template to demonstrate the various styling blocks available. By isolating \texttt{packages.tex}, \texttt{commands.tex}, and \texttt{theorem\_styles.tex}, the \texttt{main.tex} remains clean and dedicated strictly to document structure.

\subsection{Standard Core Blocks}

\begin{definition}
A \textbf{Definition block} is used to declare and explain a new mathematical concept or term. It utilizes the unified standard styling.
\end{definition}

\begin{theorem}
A \textbf{Theorem block} is intended for major educational results. It utilizes the unified standard styling exactly like Definition and Lemma.
$$ a^2 + b^2 = c^2 $$
\end{theorem}

\begin{lemma}
A \textbf{Lemma} highlights a stepping-stone result. It also utilizes the unified standard styling with a strong left vertical alignment line.
\end{lemma}

\subsection{Distinguishable Enhancements}

\begin{example}
An \textbf{Example block} demonstrates applications purely with standard formatting—no special borders or backgrounds, just clear plain text layout.
\end{example}

\begin{remark}
A \textbf{Remark block} provides supplementary information, edge cases, or commentary. It sits subtly on a light gray background with no borders.
\end{remark}

\begin{important}
An \textbf{Important block} is reserved for crucial warnings, common pitfalls, or core takeaways that the reader absolutely must not miss. It sticks out with a bold, full border enclosure.
\end{important}

\subsection{Technical Implementations}

\begin{codeexample}[language=Python]
# We use the Code Example block to distinctly encapsulate programming 
# snippets or configuration strings, utilizing a subtle gray frame 
# and clean line numbering.

def compile_latex(file):
    print(f"Compiling {file}...")
    return compile(file)
\end{codeexample}

\begin{algorithm}
The \textbf{Algorithm block} is meant to present structured pseudocode or computational sequences. It utilizes strong top and bottom bounding borders to structurally isolate it from standard literature.

Initialize $x \gets 0$ \\
Loop while $x \lt 10$ \\
\quad $x \gets x + 1$ \\
Output $x$
\end{algorithm}
